% !TEX program = xelatex
% 실험16 Gated Late Fusion — fold 별로 λ 가 게이트를 얼마나 쓰게 했는지 (PPT 슬라이드 한 장, 16:9).
% 수치 출처: outputs/gmu_nested/gated_nested/results.json
%   (E1_learn_beta_lambda, E1_learn_beta_gate 의 tab·img min/max, E1_learn_beta_inner_scores, A_primary, E1_learn_beta).
% 분류: 두 게이트(tab, img)의 시험 환자 값 중 1 에서 가장 먼 거리 d
%   꺼짐 = λ=100 (d < 0.001) / 거의 꺼짐 = d ≤ 0.03 / 조금 작동 = d > 0.2 (해당 fold 는 0.21–0.24)
\documentclass[tikz,border=0pt]{standalone}
\usepackage{fontspec}
\setmainfont{Noto Sans CJK KR}
\setsansfont{Noto Sans CJK KR}
\usepackage{amsmath}
\usetikzlibrary{calc,arrows.meta}

\definecolor{navy}{HTML}{2F3B52}
\definecolor{ink}{HTML}{1A1A1A}
\definecolor{ink2}{HTML}{5A5E66}
\definecolor{bg}{HTML}{FAFAF7}
\definecolor{off}{HTML}{8A8F98}
\definecolor{offlight}{HTML}{E4E5E8}
\definecolor{near}{HTML}{B9BDC4}
\definecolor{nearlight}{HTML}{F1F2F4}
\definecolor{on}{HTML}{D2691E}
\definecolor{onlight}{HTML}{FBE3CF}
\definecolor{good}{HTML}{3E6FB0}
\colorlet{offtext}{off!85!black}
\colorlet{neartext}{ink2}
\colorlet{ontext}{on}
\definecolor{boxyellow}{HTML}{FDEEB8}

\tikzset{every node/.style={text=ink}}

% fold 칸: \fcell{x}{y}{λ}{게이트 범위}{상태}{Δ}  상태 = off | near | on
\newcommand{\fcell}[6]{%
  \begin{scope}[shift={(#1,#2)}]
    \fill[#5light, rounded corners=5pt] (-1.75,-1.15) rectangle (1.75,1.15);
    \draw[#5, line width=1pt, rounded corners=5pt] (-1.75,-1.15) rectangle (1.75,1.15);
    \node[font=\Large\bfseries] at (0,0.5) {$\lambda=#3$};
    \node[font=\small, text=ink2] at (0,-0.12) {게이트 #4};
    \node[font=\small\bfseries, text=#5text] at (0,-0.68) {\csname lab#5\endcsname\; {\color{ink2}\mdseries ($\Delta$ #6)}};
  \end{scope}%
}
\newcommand{\laboff}{꺼짐}
\newcommand{\labnear}{거의 꺼짐}
\newcommand{\labon}{조금 작동}

\begin{document}
\begin{tikzpicture}[x=1cm,y=1cm]
\fill[bg] (0,0) rectangle (25.4,14.29);

% ---------------- 제목 ----------------
\node[anchor=west, font=\huge\bfseries] at (0.7,13.35) {이번 실험에서 $\lambda$가 게이트를 어떻게 조절했나};
\node[anchor=west, font=\normalsize, text=ink2] at (0.75,12.5)
  {$\lambda$ = 게이트를 얼마나 쓰게 할지 정하는 값. fold마다 학습 환자 안에서 성능이 가장 좋은 값을 자동으로 고름};

% ---------------- λ 눈금 ----------------
\begin{scope}[shift={(0,10.75)}]
  \shade[left color=on!75, right color=off!80, rounded corners=4pt] (6.2,-0.17) rectangle (19.2,0.17);
  \foreach \i/\v in {0/0.01, 1/0.1, 2/1, 3/10, 4/100} {
    \fill[white] ({6.7+\i*3.0},0) circle (0.11);
    \draw[navy, line width=0.6pt] ({6.7+\i*3.0},0) circle (0.11);
    \node[font=\normalsize\bfseries] at ({6.7+\i*3.0},0.62) {\v};
  }
  \node[anchor=east, align=right, font=\normalsize\bfseries, text=on] at (5.9,0.12) {$\lambda$ 작음};
  \node[anchor=east, align=right, font=\small, text=ink2] at (5.9,-0.5) {환자마다 게이트 값이 달라짐};
  \node[anchor=west, align=left, font=\normalsize\bfseries, text=offtext] at (19.5,0.12) {$\lambda$ 큼};
  \node[anchor=west, align=left, font=\small, text=ink2] at (19.5,-0.5) {게이트 꺼짐 (모두 1)};
  \node[anchor=west, align=left, font=\small, text=ink2] at (19.5,-0.95) {= 기존 late fusion과 같음};
\end{scope}

% ---------------- fold 격자 ----------------
\foreach \i in {1,...,5} \node[font=\normalsize\bfseries, text=ink2] at ({5.85+(\i-1)*3.85},8.75) {fold \i};
\node[anchor=east, font=\Large\bfseries] at (3.85,7.35) {OS};
\node[anchor=east, font=\Large\bfseries] at (3.85,4.85) {PFS};

\fcell{5.85}{7.35}{1}{0.99--1.01}{near}{0.000}
\fcell{9.70}{7.35}{100}{1.00}{off}{0.000}
\fcell{13.55}{7.35}{0.1}{0.76--1.15}{on}{+0.001}
\fcell{17.40}{7.35}{1}{0.97--1.02}{near}{$-$0.001}
\fcell{21.25}{7.35}{0.1}{0.77--1.16}{on}{+0.007}

\fcell{5.85}{4.85}{1}{0.99--1.02}{near}{+0.001}
\fcell{9.70}{4.85}{100}{1.00}{off}{0.000}
\fcell{13.55}{4.85}{0.1}{0.79--1.09}{on}{$-$0.001}
\fcell{17.40}{4.85}{100}{1.00}{off}{0.000}
\fcell{21.25}{4.85}{1}{0.98--1.01}{near}{$-$0.001}

\node[anchor=west, font=\scriptsize, text=ink2] at (4.1,3.35)
  {게이트 = 시험 환자의 임상·영상 게이트 값 범위 (1이면 late fusion과 같음) \quad $\Delta$ = 이 fold에서 late fusion 대비 C-index 차이};

% ---------------- 아래 왼쪽: 개수 ----------------
\begin{scope}[shift={(0.7,0.45)}]
  \fill[white, rounded corners=6pt] (0,0) rectangle (11.6,2.55);
  \draw[navy!60, line width=0.7pt, rounded corners=6pt] (0,0) rectangle (11.6,2.55);
  \node[anchor=west, font=\large\bfseries] at (0.35,2.0) {10개 fold 중};
  \fill[off]  (0.4,0.45) rectangle (3.4,1.25);
  \fill[near] (3.4,0.45) rectangle (7.4,1.25);
  \fill[on]   (7.4,0.45) rectangle (10.4,1.25);
  \node[font=\normalsize\bfseries, text=white] at (1.9,0.85) {꺼짐 3};
  \node[font=\normalsize\bfseries, text=ink]   at (5.4,0.85) {거의 꺼짐 4};
  \node[font=\normalsize\bfseries, text=white] at (8.9,0.85) {조금 작동 3};
  \node[anchor=west, font=\small, text=ink2] at (4.1,2.0) {작동한 fold도 차이는 $-$0.001 \textasciitilde{} +0.007};
\end{scope}

% ---------------- 아래 오른쪽: 왜 / 결론 ----------------
\begin{scope}[shift={(12.7,0.45)}]
  \fill[boxyellow!70, rounded corners=6pt] (0,0) rectangle (12.0,2.55);
  \draw[navy!60, line width=0.7pt, rounded corners=6pt] (0,0) rectangle (12.0,2.55);
  \node[anchor=north west, text width=11.3cm, align=left, font=\normalsize, inner sep=0pt] at (0.35,2.25)
    {\textbf{왜?} $\lambda$ 후보 5개 성능 차이가 fold마다 0.01 미만 → 써도 안 써도 비슷\\[4pt]
     \textbf{결과} 게이트로 얻은 이득이 거의 없음 → late fusion과 같은 성능\\
     {\small\color{ink2}\phantom{\textbf{결과} }OS 0.725 vs 0.723, PFS 0.674 vs 0.674}};
\end{scope}

\end{tikzpicture}
\end{document}
